% 第10批：《大学物理简明教程（第5版）》第10章“波动光学”
% 原书及PDF第270—271页，选择题10.1.1—10.1.18。
% 独立题意摘要、详细原理与核验；图形题保留判断边界，待独立二校。
\subsection{选择题：干涉、衍射、光栅和偏振}

\sourceproblem{10.1.1}{270}{270}
\problemtext 同一单色光在空气和玻璃中传播相同时间，比较两条光路的几何长度与光程。
\workings
折射率为$n$时速度为$c/n$，相同时间$t$的几何长度
$s=ct/n$，故空气中路程更长。其光程$ns=ct$，两介质中相等。
\[
\boxed{s_{\rm air}\ne s_{\rm glass},\quad
n_{\rm air}s_{\rm air}=n_{\rm glass}s_{\rm glass}=ct}.
\]
原题选\textbf{C}。\crosscheck 不可把几何路程与光程当成同一概念；题目明确比较的是\emph{相同时间}。

\sourceproblem{10.1.2}{270}{270}
\problemtext 杨氏双缝实验中，希望使屏上干涉条纹间距增大，比较改变屏距、双缝间距、缝宽和波长的做法。
\workings
旁轴条件下条纹间隔
$\Delta x=\lambda D/d$。保持光源和屏距，减小双缝间距$d$便使条纹变疏，
故原题选\textbf{B}。
\crosscheck 单独缩小每条缝的宽度主要改变衍射包络，不是双缝中心条纹间距。

\sourceproblem{10.1.3}{270}{270}
\problemtext 双缝间距$0.50\,\mathrm{mm}$、波长$600\,\mathrm{nm}$，屏距$1.20\,\mathrm m$；求相邻明条纹距离。
\workings
\[
\boxed{\Delta x=\frac{\lambda D}{d}
=\frac{600\times10^{-9}(1.20)}{0.50\times10^{-3}}
=1.44\,\mathrm{mm}}.
\]
原题选\textbf{D}。\crosscheck 计算前须将mm与nm统一换为m；其结果不依赖第几级明纹。

\sourceproblem{10.1.4}{270}{270}
\problemtext 两块玻璃板构成左缘接触、向右增厚的空气劈尖，现上玻璃板绕左接触棱作微小顺时针转动，判断等厚条纹间距及平移方向。
\workings
劈尖斜率$\alpha$减小，反射等厚干涉相邻条纹间隔
$\Delta x\simeq\lambda/(2\alpha)$变大。
固定干涉级次的厚度$t_m$保持不变，
$x_m=t_m/\alpha$增大，故条纹远离左端棱边。
原题选\textbf{B}。
\crosscheck 判断方向须固定题目指定的\emph{左缘为接触棱}；
若转动方向相反，两个结论都可能反向。

\sourceproblem{10.1.5}{270}{270}
\problemtext 波长$\lambda$的单色光垂直入射夹在空气与高折射率衬底之间的透明薄膜，膜折射率为$n$且$1<n<n_{\rm sub}$。求使两束主要反射光相长干涉的最小非零膜厚。
\workings
空气到膜及膜到衬底均是\emph{疏到密}反射，两束反射都发生$\pi$相位突变，因而相互比较无附加半波差。
相长条件$2nt=m\lambda$；取$m=1$得
\[
\boxed{t_{\min}=\frac{\lambda}{2n}}.
\]
原题选\textbf{D}。\crosscheck 若膜折射率大于衬底则仅有一次半波损失，结果将不同；须保留本题的折射率大小顺序。

\sourceproblem{10.1.6}{270}{270}
\problemtext 牛顿环实验用液体替代空气后，第10条亮环直径由$1.42$ cm变为$1.28$ cm。求液体折射率。
\workings
对于固定级次与曲率半径，环半径平方与介质折射率成反比，
$n=(D_{\rm air}/D_{\rm liq})^2$，因此
\[
\boxed{n=(1.42/1.28)^2\approx1.231}.
\]
原题选\textbf{C}。\crosscheck 不可只取直径之比$1.109$；光程差由介质折射率进入，环半径为平方根关系。

\sourceproblem{10.1.7}{270}{270}
\problemtext 迈克耳孙干涉仪一臂放入长度$t=0.50$ cm的气体小室，波长$500$ nm；小室由真空缓慢充气后有60条干涉条纹移过视场，求气体折射率。
\workings
干涉仪光线往返气室，额外光程为$2(n-1)t$。移过$N=60$条纹意味着$2(n-1)t=N\lambda$，故
\[
\boxed{n=1+\frac{60(500\times10^{-9})}{2(0.005)}
=1.003}.
\]
原题选\textbf{C}。\crosscheck 迈克耳孙一臂\emph{往返}，缺少因子2会误得1.006。

\sourceproblem{10.1.8}{270--271}{270--271}
\problemtext 迈克耳孙干涉仪的一个光臂插入折射率$n$、厚度$d$的透明薄片，求该臂单程光程增加量。
\workings
用光学介质替换同厚空气，单程光程由$d$变为$nd$，
故
\[
\boxed{\delta_{\rm one-way}=(n-1)d}.
\]
原题选\textbf{E}。\crosscheck 题问的是该\emph{单程光路}光程改变；
若问干涉仪双程的总光程差则为$2(n-1)d$，不能混淆。

\sourceproblem{10.1.9}{271}{271}
\problemtext 迈克耳孙一臂加一透明膜，使两束光的光程差改变恰好一个真空波长$\lambda$，求膜厚。
\workings
往返穿过膜两次，干涉两束光程差改变量为$2(n-1)d$，
故
\[
\boxed{d=\frac{\lambda}{2(n-1)}}.
\]
原题选\textbf{D}。\crosscheck 第10.1.8题的单程光程增加与本题\emph{往返干涉}要区分，不能直接把前题结果作为总光程差。

\sourceproblem{10.1.10}{271}{271}
\problemtext 单缝夫琅禾费衍射装置中，在中央明纹位置不变时缩小狭缝宽度$a$，判断衍射条纹的疏密。
\workings
单缝暗纹满足$a\sin\theta_m=m\lambda$。
固定$m,\lambda$，缩小$a$使$|\theta_m|$增大，条纹向外展开，变\textbf{疏}。
原题选\textbf{B}。\crosscheck 不能用双缝间距公式替代单缝衍射暗纹条件。

\sourceproblem{10.1.11}{271}{271}
\problemtext 波长$\lambda$垂直入射单缝，观察到$\theta=30^\circ$方向波面可分成3个半波带，求缝宽$a$。
\workings
半波带数量$N=2a\sin\theta/\lambda$，给$N=3$，故
\[
\boxed{a=\frac{3\lambda}{2\sin30^\circ}=3\lambda}.
\]
原题选\textbf{C}。\crosscheck 3个\emph{半波带}代表光程差$3\lambda/2$而非$3\lambda$。

\sourceproblem{10.1.12}{271}{271}
\problemtext 波长$500$ nm单色光经宽$a=0.25$ mm的单缝及凸透镜，在焦面观察衍射；两侧第3条暗纹间隔为12 mm，求透镜焦距。
\workings
第$m$级暗纹坐标近似$x_m=mf\lambda/a$，
两侧$m=3$暗纹间隔$d=6f\lambda/a$。
\[
\boxed{f=\frac{da}{6\lambda}
=\frac{0.012(0.25\times10^{-3})}{6(500\times10^{-9})}
=1.00\,\mathrm m}.
\]
原题选\textbf{B}。\crosscheck 两侧第3级暗纹相距$6$个单位位移，非$3$个。

\sourceproblem{10.1.13}{271}{271}
\problemtext 由$N$条相同透光缝构成的理想光栅，比较其中央主极大的强度与单条缝的中央衍射强度。
\workings
中央方向相干叠加振幅为单缝的$N$倍，强度为其平方，
\[
\boxed{I_{\max,N}=N^2I_{\max,1}}.
\]
原题选\textbf{D}。\crosscheck 这里比较单色、等振幅、相干的\emph{峰强度}，而不是总透射能量。

\sourceproblem{10.1.14}{271}{271}
\problemtext 光栅常量为$d=a+b$，垂直入射波长$\lambda$光，求第$k$级主极大的衍射角$\varphi$满足的方程。
\workings
相邻缝程差为$d\sin\varphi$，主极大同相要求
\[
\boxed{d\sin\varphi=k\lambda},\qquad k\in\mathbb Z.
\]
原题选\textbf{D}。\crosscheck 狭缝宽度$a$还决定单缝衍射包络，不能把它代入主极大光栅常量$d$。

\sourceproblem{10.1.15}{271}{271}
\problemtext 光栅中所有偶数级主极大恰与单缝衍射暗纹重合而缺级，求透光缝宽$a$与不透明部分宽$b$的关系。
\workings
单缝零强度条件$a\sin\theta=m\lambda$，
光栅主极大条件$(a+b)\sin\theta=k\lambda$。
若每个偶数级$k=2m$消失，须
$(a+b)/a=2$，即
\[
\boxed{a=b}.
\]
原题选\textbf{B}。\crosscheck 严格说偶数级消失的是由单缝包络造成的\emph{缺级}，而不是相干叠加的零级主极大。

\sourceproblem{10.1.16}{271}{271}
\problemtext 强度$I_0$的自然光顺序穿过两偏振片，透振方向夹角$45^\circ$，求出射强度。
\workings
自然光经过理想首偏振片强度减半，第二片按马吕斯定律，
\[
\boxed{I=\frac{I_0}{2}\cos^245^\circ=\frac{I_0}{4}}.
\]
原题选\textbf{B}。\crosscheck 不能从自然光直接写$I_0\cos^2\theta$，忽略第一片的$1/2$。

\sourceproblem{10.1.17}{271}{271}
\problemtext 自然光从空气按布儒斯特角入射玻璃面，判定反射光的偏振状态。
\workings
布儒斯特角满足$\tan i_B=n_2/n_1$，此时反射光的平行入射面$p$分量为零，只剩\emph{垂直于入射面}的$s$分量，
故反射光是完全线偏振光。
原题选\textbf{C}。\crosscheck 反射光的偏振方向不是\emph{平行}入射面；完全偏振的是反射光，不是所有折射光。

\sourceproblem{10.1.18}{271}{271}
\problemtext 自然光按布儒斯特角射向平板玻璃上表面，再按原图追踪内部反射后从玻璃板射出的标号2的反射相关光束，判断其偏振状态。
\workings
布儒斯特入射时首次反射的$p$分量消失；
平行两表面满足折射角$r=90^\circ-i_B$，玻璃向空气反射也恰达到反向布儒斯特条件。
因此原图标号2对应的二次\emph{反射}光只具有垂直入射面的$s$分量，是线偏振光，
\[
\boxed{\bm E_2\perp\text{入射面}}.
\]
原题选\textbf{C}。
\crosscheck \textbf{图题待二校：}须把标号2识别为经内部表面布儒斯特反射的光线，而不是一次透射光；若图中光线路径有别，偏振结论也须相应调整。
